% !TEX program = lualatex
\documentclass[a4paper,11pt]{ltjsarticle}

\usepackage{amsmath,amssymb}
\usepackage{booktabs}
\usepackage{geometry}
\usepackage{hyperref}
\geometry{margin=25mm}

\title{式(2)におけるGPプラント表現と\\確率的ゆらぎの説明}
\author{}
\date{\today}

\newcommand{\GP}{\mathrm{GP}}

\begin{document}
\maketitle

\section{式(2)の位置づけ}

レジュメにおける式(2)は，本研究で多目的GPを「プラント」として見たときの
集団遷移を表す式である．
式は次のように書かれる．
\begin{equation}
P_{g_{\ell+1}}
=
F(P_{g_\ell},u_\ell,\xi_\ell).
\end{equation}

これは，制御ステップ \(\ell\) において，現在の集団 \(P_{g_\ell}\) に対して
BOが決定した制御入力 \(u_\ell\) を与え，GPを \(k_\ell\) 世代進化させた結果，
次の制御更新時点の集団 \(P_{g_{\ell+1}}\) が得られる，という意味である．

直感的には，
\[
\text{次の集団}
=
\text{GPの進化処理}
(\text{現在の集団},\text{制御入力},\text{乱数の影響})
\]
を1つの式で表している．

\section{各記号の意味}

\begin{table}[htbp]
\centering
\caption{式(2)に含まれる記号の意味}
\begin{tabular}{ll}
\toprule
記号 & 意味 \\
\midrule
\(g\) & GPの世代番号 \\
\(\ell\) & BOによる制御更新ステップ番号 \\
\(P_{g_\ell}\) & 世代 \(g_\ell\) における現在集団 \\
\(P_{g_{\ell+1}}\) & 次の制御更新時点 \(g_{\ell+1}\) における集団 \\
\(u_\ell\) & BOが決定する制御入力 \\
\(u_\ell=(p_{c,\ell},p_{m,\ell},k_\ell)\) & 交叉率，突然変異率，更新周期 \\
\(F\) & GPの進化作用全体 \\
\(\xi_\ell\) & GP進化中に生じる確率的ゆらぎ \\
\bottomrule
\end{tabular}
\end{table}

ここで \(F\) は，親選択，交叉，突然変異，複製，評価，
NSGA-IIによる次世代選択などをまとめたものである．
したがって，式(2)は個々の操作を細かく展開する式ではなく，
GP全体を制御対象として抽象化した式である．

\section{\texorpdfstring{確率的ゆらぎ \(\xi_\ell\) とは何か}{確率的ゆらぎ xi ell とは何か}}

\(\xi_\ell\) は，GPの進化過程に含まれる乱数の影響をまとめた記号である．
GPでは，同じ集団から始め，同じ交叉率 \(p_c\)，突然変異率 \(p_m\)，
更新周期 \(k\) を用いたとしても，乱数の出方が異なれば次に得られる集団は完全には一致しない．
この「同じ条件でも結果が少し変わる要因」をまとめて
確率的ゆらぎと呼ぶ．

具体的には，\(\xi_\ell\) には次のような要素が含まれる．

\begin{itemize}
  \item トーナメント選択などで，どの個体が親として選ばれるか．
  \item どの個体同士が交叉するか．
  \item プログラム木のどの位置で交叉が起こるか．
  \item 突然変異でどのノードや部分木が選ばれるか．
  \item 突然変異によってどのような新しい部分木が生成されるか．
  \item 複製操作でどの個体がそのまま次世代候補に残るか．
  \item 同順位の個体が存在する場合に，どの個体が選ばれるか．
\end{itemize}

つまり，\(\xi_\ell\) はBOが直接決める値ではない．
BOが決めるのは
\[
u_\ell=(p_{c,\ell},p_{m,\ell},k_\ell)
\]
であり，\(\xi_\ell\) はその制御入力のもとでGPを実行したときに自然に発生する
確率的なばらつきである．

\section{なぜこの記号を入れるのか}

\(\xi_\ell\) を式に含める理由は，GPの応答が完全には決定的でないことを明示するためである．
仮に同じ \(P_{g_\ell}\) と同じ \(u_\ell\) を与えても，
乱数系列が異なれば \(P_{g_{\ell+1}}\) や報酬 \(r_\ell\) は変わり得る．
したがって，BOが観測する報酬は，
制御入力だけで一意に決まる完全な値ではなく，
確率的なばらつきを含んだ観測値である．

本研究では，この性質を含めたうえで，
BOが
\[
r=f(x,p_c,p_m,k)
\]
を学習し，現在の状態 \(x_\ell\) に応じて次の制御入力を選択する．
すなわち，\(\xi_\ell\) は
「BOが制御する部分」と「GP実行時に不可避に生じるランダムな部分」を
区別するための記号である．

\section{発表時の説明例}

発表では，次のように説明すると分かりやすい．

\begin{quote}
式(2)は，GPを制御対象，つまりプラントとして見たときの集団の変化を表しています．
\(P_{g_\ell}\) は現在の集団，\(u_\ell\) はBOが決める
交叉率・突然変異率・更新周期です．
ただし，GPには親選択，交叉点，突然変異点など乱数で決まる部分があるため，
同じ操作率を使っても次の集団が完全に同じになるとは限りません．
この乱数による影響をまとめて \(\xi_\ell\) と書いています．
つまり，BOは \(u_\ell\) を制御しますが，
\(\xi_\ell\) によるばらつきを含んだ結果を観測しながら，
次に良さそうな操作率と更新周期を学習していく，という構造です．
\end{quote}

\section{要点}

\begin{itemize}
  \item 式(2)は，GPの世代進化を制御対象としてまとめたプラント表現である．
  \item \(u_\ell=(p_c,p_m,k)\) はBOが決める制御入力である．
  \item \(\xi_\ell\) は，GP内部の乱数による確率的なばらつきを表す．
  \item \(\xi_\ell\) を入れることで，同じ制御入力でも結果が完全には固定されないことを明示できる．
  \item BOは，このばらつきを含む観測結果から状態依存の制御方策を学習する．
\end{itemize}

\end{document}
